% ================================================================
%  chapters/parts/ch10-c.tex
%  第 10 章  多相催化反应动力学（分片 c）
%  源：S2（下册）第 10 章，书内 p.779—806 / PDF p.268—295
%  覆盖：10.5 动力学测定方法和实验装置；10.6 非稳态催化过程动力学；10.7 结论
%  说明：本分片负责 10.5 与 10.6（分片 b 负责 10.3 与 10.4）。
%        10.5.1—10.5.4 的源文本区间为书内 p.779—792 / PDF p.268—281。
% ================================================================

\subsection{动力学测定方法和实验装置}
\label{sec:ch10-measure}

在一般多相催化反应器上得到的实验结果里，化学反应与传递过程总是不同程度地耦合在一起。
用这种耦合数据，无论设计反应器还是改进催化剂，都是困难的。因此必须先靠适当的实验工具
和条件把两者分离开，才谈得上求取\textbf{本征动力学参数}（速率常数、活化能、反应级数、
吸附平衡常数）与\textbf{传递过程参数}（有效扩散系数、传质与传热系数）。分离的手段有两类：
选对反应器与操作区，以及用非稳态实验给稳态实验补上它拿不到的信息。
\src{S2 §10.5 引言，书内 p.779 / PDF p.268}

\begin{probout}
本章的前半部分交付速率方程，而速率方程能否被信任，取决于数据是在哪个\textbf{动力学区}
测出来的。本分片回答四个问题：（1）实验室反应器种类繁多，凭什么从中选一个——
它能测什么、不能测什么；（2）「已排除外扩散」「已排除内扩散」不是一句声明，而是
可以执行的\textbf{三组实验}与三个\textbf{定量判据}，它们具体怎么做、阈值是多少；
（3）稳态动力学给出的是速率对流体相组成的依赖，它\textbf{给不出吸附态的浓度与寿命}，
这个缺口由哪些非稳态方法补、各自的时间尺度是多少；（4）催化剂活性随时间变化、
多稳态与振荡这些「非稳态」现象，如何做动力学处理。
\end{probout}

\subsubsection{实验室反应器的选型逻辑}

反应器种类繁多，\textbf{但并不是都可以用来求取动力学参数}。选型依据是催化反应的特点
和所要获取的信息，其中有三个条件优先于其他一切考虑\src{S2 §10.5.1，书内 p.779 /
PDF p.268}：

\begin{keycond}
（1）\textbf{恒温性}——温度对反应速率的影响是指数性的，反应器内任何温度梯度都会直接
变成速率误差；对复杂反应，温度偏差还会改变选择性。（2）\textbf{停留（接触）时间的确切性
或均一性}——数据要与模型比较，就必须知道每个流体元在催化剂上待了多久。（3）\textbf{产物
取样和分析是否容易}——分析精度往往才是真正的限制环节。至于反应器制作是否价廉简易，
\textbf{应服从实验需要}，不能作为选型的首要依据。
\end{keycond}

\textbf{（1）积分反应器。}管式反应器中装填足量催化剂以达到较大转化率，进出口组成差别
显著，无法用某个平均值代表全反应器的组成。因此积分反应器只能给出\textbf{转化率（或
生成率）对空时的积分数据}，但在分析技术和方法上要求不高。它在动力学研究中又分两种。
\textbf{恒温积分反应器}简单价廉，只要可能总应首先考虑；其困难在于转化率高、反应热稍大
就难以恒温。常用对策有三条：\textbf{减小管径}（使径向温度均匀，管径为催化剂粒径 4 倍
以上时减小管径对恒温性的改善仍很明显，但管径过小会加剧沟流引起的壁效应，使转化率
偏低）、\textbf{用恒温导热介质}（熔融金属、熔盐、流沙浴；流沙浴可达 $1000$\,\degC，
控温快而简易，即使放热量较大的氧化反应也可用）、\textbf{用惰性物质稀释催化剂}
（强放热反应用刚玉、石英砂等大热容粒子避免热点）。进一步可采用\textbf{非等比稀释}：
入口处加大稀释比、随转化加深线性降低稀释比，使轴向温度梯度接近于零。
\textbf{绝热积分反应器}为直径均一、装填均匀、绝热良好的圆管；物料预热到一定温度后
通入，沿轴向可测出与反应热量和动力学规律相应的温度分布 $T$–$W$ 曲线。对这条曲线作
解析微分，得到一系列 $\dd T/\dd W$ 与 $\dd^{2}T/\dd W^{2}$；若催化剂有效导热率 $\lambda_e$
已知，且径向梯度与气固间梯度小到可忽略，把这些导数代入反应管的热量衡算（轴向对流项、
有效导热项与反应热源项三项平衡），即可解出不同温度下的 $r_1$；再由能量衡算换算成同一
温度下各组分的摩尔分数 $Y_i$。有了一套不同入口温度下的 $(r_i,\,Y_i)$ 数据，动力学参数
就可以进一步求得——\textbf{但只能用于简单反应}。\src{S2 §10.5.1，书内 p.779—780 /
PDF p.268—269}

\textbf{（2）微分反应器。}若通过床层的转化率很低，进出口组成差别小得足以用平均值代表
全床层组成，却又足够用某种分析方法确定浓度差，则 $\Delta c/\Delta t$ 可近似为
$\dd c/\dd t$ 并等于反应速率 $r$，从而得到与平均组成对应的\textbf{点速率}。单程流通的
管式微分反应器中转化率一般应低于 $5\%$（个别情况允许到 $10\%$）。优点是转化低、发热少、
易恒温、构制简单；但有两个严重问题：一是所得数据常是\textbf{初速}，很难拼出一个合成的
进料来模拟复杂反应在高转化下的组成（补救办法是在微分反应器前串接一个积分转化器，为它
提供所需进料）；二是\textbf{分析精度要求很高}，大多数方法满足不了。\src{S2 §10.5.1，
书内 p.780—781 / PDF p.269—270}

\textbf{（3）循环反应器（无梯度反应器）。}积分反应器分析容易但不能确保理想行为，微分
反应器能确保理想行为但分析困难，\textbf{循环反应器同时具备两者的优点}。反应流体的大部分
用泵在回路中循环，同时连续引入一小股新鲜进料 $F_0$，并从出口放出一股流出物以保持系统
恒压。流体可外循环，也可内循环。设循环量 $F_R$，进料中反应组分 B 的摩尔分数为 $y_0$，
入床层前为 $y_{\mathrm{in}}$，出口为 $y_f$，则按物料衡算

\begin{equation}
x=y_f-y_{\mathrm{in}}=\frac{y_f-y_0}{1+F_R/F_0}
\label{eq:recycle-x}
\end{equation}

当 $F_R\gg F_0$ 时 $y_{\mathrm{in}}\to y_f$、$y_f-y_{\mathrm{in}}\to 0$，反应速率由
微分形式给出

\begin{equation}
r=\frac{\dd x}{\dd W/(F_0+F_R)}\approx\frac{y_f-y_0}{W/F_0}
\label{eq:recycle-r}
\end{equation}

$F_R/F_0$ 一般为 $20$—$40$，远大于 1，相当于把 $y_f-y_{\mathrm{in}}$ 这个微小差值
放大成分析上容易测准的 $y_f-y_0$。由于通过床层的转化率很低，床层温度变化很小；又由于
通过床层的流体量相当大、线速大，\textbf{外扩散影响也可以消除}。代价从分析精度转移到
循环泵：泵不能沾污反应混合物，滞留量要小，循环量要大（$4$~L\,min$^{-1}$ 以上）。
外循环反应器另有两条不足：循环气须冷却到泵体能忍受的温度、出泵后再预热到反应温度；
自由体积与催化剂体积之比高达 $10$—$100$，时间常数很大，从一个操作条件换到另一个
需要较长时间才达稳态，且有利于副反应进行。\textbf{内循环反应器}把循环系统与反应系统
合为一体，自由体积降为催化剂体积的 $3$—$5$ 倍，时间常数显著缩短；装置上部为催化床层
和风道、下部为涡轮叶片，高速转动时强制气流通过催化床再经风道吸回，循环气量可达
$150$~L\,min$^{-1}$，浓度和温度更易达到无梯度状态。它要求轴转速达每分钟数千转，并能
在高温高压下无泄漏、无沾污；目前特殊轴承材料可在 $590$\,\degC 及 $15$~MPa 下操作，
采用磁驱动、全玻璃（或石英）的设计还能避免金属结构材料对催化反应的干扰，且结构简单、
价格低廉。\src{S2 §10.5.1，书内 p.781—782 / PDF p.270—271}

\textbf{（4）搅拌反应器（CSTR）。}当循环比足够高时，循环反应器实际上就是一种全返混
反应器。实验室中与 CSTR 相当的是\textbf{旋转篮式反应器}：催化剂颗粒装在金属网编织的
筐篮中（一般整齐排列），筐篮连于转轴，在反应容器中高速回转，流入的流体在瞬时内与容器
内原有流体完全均匀混合，\textbf{组成与流出物完全相同}。最早的是 Carberry 提出的
notre dame 式；后来的 ICI 式提高了转速、改用磁驱动以防沾污和泄漏并缩小了自由体积；
再后又用加挡板和固定筐篮、使容器高速旋转的办法改进测温和加大流体穿透催化剂的相对速度。
另一种产生气体搅拌效果的办法是让器内气体往复振荡（外部交变电磁场驱动活塞，或由蜂鸣器
带动膜片）。\src{S2 §10.5.1，书内 p.782—783 / PDF p.271—272}

\textbf{（5）流化床反应器与微反应器。}如果催化剂活性在反应过程中逐步上升（诱导期）
或下降（失活），就应改用流化床或微反应器：\textbf{前者适于失活时间以秒计、且有较多
催化剂量可供应用的情况；后者宜于活性变化稍慢而催化剂供应量很少的情况}。流化床反应器
曾用于邻二甲苯氧化动力学研究（内径 $17$~mm、高 $8.2$~m 的不锈钢竖管，含 SO$_2$ 的空气
带入细粒氧化钒催化剂，达稳定两相流动后通入定量邻二甲苯，几分钟即达定态，从上、中、下
三个取样点经烧结不锈钢过滤器滤去催化剂后取样分析）。\src{S2 §10.5.1，书内 p.783 /
PDF p.272}

\textbf{（6）无梯度性能的比较与检验。}动力学反应器基本上就是指无梯度反应器（循环式与
搅拌式）。涡轮或筐篮的转速（或循环量）对无梯度状况有显著影响，后者必须用两类实验检验：
\textbf{用传质实验检验外扩散阻力}（把一定尺寸的萘锭与催化剂装在一起，在 $30$\,\degC 下
测定氮气流中的萘浓度，求出传质因子 $j_D$），以及\textbf{用返混实验表征混合均匀度}
（在进器气流中注入示踪气脉冲，测出器气流中示踪物浓度随时间的变化）。无梯度反应器的
性能比较见表~\ref{tab:ch10-reactor}。\src{S2 §10.5.1，书内 p.783—784 / PDF p.272—273}

选型不是「哪个更好」，而是\textbf{先按所需信息过滤，再按恒温性、接触时间均一性、
分析难度排序}：要宽转化率范围的积分数据选积分床；要点速率选微分床或无梯度反应器；
体系活性快变选流化床或脉冲微反；体系只允许很少催化剂时选微反。表~\ref{tab:ch10-reactor}
把这一逻辑与「能测／不能测」并列，是本节选型的总表。

\begin{table}[tbp]
\centering
\caption{实验室动力学反应器的选型对照。先看第三列\textbf{不能测什么}——反应器的用途
边界比它的优点更需要提前确认；第四列给出该类型成立所必需的约束与本书给出的经验数值。
「流化床」一栏兼收本书讨论的气固流化床，液固浆态体系据本书给出的液固床层条件与
常规实践补充（见本节末注）。\src{据 S2 §10.5.1 与表 10-3，书内 p.779—784 / PDF
p.268—273 整理。}}
\label{tab:ch10-reactor}
{\srcsize
\begin{tabularx}{\twocolwidth}{@{}p{17mm}p{40mm}p{40mm}X@{}}
\toprule
反应器 & 能测什么 & 不能测什么 & 关键约束（恒温性／时间均一性／分析／解析／成本） \\
\midrule
积分固定床（恒温） & 转化率（或生成率）对空时的\textbf{积分数据}；宽转化率范围内的表观
速率；总体速率方程的积分检验 & 给不出某一点的点速率——进出口组成差别大，任何平均值都
无意义；高转化下的数据难以拆回初速 & 恒温性最差：转化高、反应热稍大即难恒温。对策为
减小管径（$\ge 4d_p$ 仍有效，过小则沟流）、恒温导热介质、惰性稀释、非等比稀释。
分析要求低，易、廉（恒温性评为「差」） \\
积分固定床（绝热） & 轴向温度分布 $T$–$W$ 曲线；解析微分得 $\dd T/\dd W$ 与
$\dd^{2}T/\dd W^{2}$，配合 $\lambda_e$ 求不同温度下的 $r_1$ 与 $Y_i$ & \textbf{只能用于
简单反应}；温度与转化率沿床层同时变化，给不出等温速率 & 需绝热良好、直径均一、装填
均匀；须假定径向梯度与气固间梯度可忽略。解析难，成本中等 \\
微分固定床（单程） & $\Delta c/\Delta t\approx\dd c/\dd t=r$；与平均组成对应的
\textbf{点速率}；易恒温、构制简单 & 数据常为\textbf{初速}，难以模拟复杂反应在高转化下的
组成；\textbf{分析精度要求极高}，多数方法不满足；接触时间均一性只算中等 & 单程转化率应
低于 $5\%$（个别允许 $10\%$）；可在其前串接积分转化器以提供所需进料。分析难、解析良、
易且廉 \\
循环反应器（外循环／内循环，无梯度） & 点速率 $r\approx(y_f-y_0)/(W/F_0)$；转化率低、
床层温差小，且线速大，\textbf{外扩散易于消除}；分析容易 & 循环泵制造难；外循环自由体积／
催化剂体积达 $10$—$100$，时间常数大，换条件后久难达稳态且利于副反应；快速失活体系
不适用 & $F_R/F_0$ 取 $20$—$40$；循环量 $>4$~L\,min$^{-1}$。内循环自由体积降为
$3$—$5$ 倍、循环量可达 $150$~L\,min$^{-1}$、可耐 $590$\,\degC 与 $15$~MPa；磁驱动
全玻璃／石英可避免金属沾污。成本：内循环难且贵、外循环中等 \\
旋转篮式 CSTR（Carberry／ICI 型） & 全返混条件下的点速率；器内组成等于出口组成，
\textbf{取样分析容易}；可测活性在变化中的催化剂 & 高速回转下的密封与测温困难；
\textbf{与管式反应器结果可能不一致}（见后文误判框）；接触时间均一性仅「良」 &
需挡板与磁驱动以缩小自由体积、防沾污；须用萘锭传质实验检验外扩散、用脉冲响应检验
返混。解析良、制作难且贵 \\
流化床（气固）／浆态（液固） & 活性在反应中变化（诱导期、失活）的动力学；液固体系
可在恒温与近似等浓度下取点速率 & 接触时间不均一、返混强；气固接触效率低；磨耗与夹带
使实际催化剂量不确定；取样须过滤 & 适合失活时间以秒计的体系；液-固床层长度应达粒径的
几百倍以上。温度均一性仅「中」，解析良，制作难 \\
脉冲微反（含 TAP 微反应器） & 吸附、脱附、扩散参数；反应网络与机理证据；\textbf{活性
变化很快而催化剂量很少}的体系；仅 Knudsen 扩散下的本征相互作用 & 给不出稳态速率方程与
高转化率下的选择性；脉冲之间表面状态被强制扰动，不代表稳态覆盖度；不能做物料衡算式的
积分检验 & TAP：脉冲宽度约 $10^{-4}$~s、弛豫时间 $10^{-4}$—$10^{3}$~s；脉冲分子数
$10^{10}$—$10^{15}$ 个，只占表面活性位数量的百万分之一；对流传质消失、仅有 Knudsen
扩散；电离室距反应器出口 $\le 1$~cm，背景压力 $<10^{-6}$~Pa \\
\bottomrule
\end{tabularx}}
\end{table}

\noindent\textbf{表~\ref{tab:ch10-reactor} 的一处边界：}「浆态（液固）」一行超出原书
§10.5.1 的展开范围。原书在此处讨论的是气固流化床，对液固体系只给出床层几何条件
（液-固床层长度应为粒径的几百倍以上）与恒温性要求。
\tocheck{未对照原页面图像：浆态反应器（间歇釜、浆态床）在本书 §10.5.1 中是否另有条目，
据 ch10-b 分片源文本（书内 p.779—784）未见；此行为按常规实践与本书液固条件补充，供回查。}

\begin{measure}
\textbf{无梯度性的两项检验（在正式取数之前完成）。}（1）\textbf{传质实验}：把一定尺寸
（如长径都为 $3$~mm 的圆柱体）的萘锭与催化剂一起装入，在 $30$\,\degC 下测通过氮气流中
的萘浓度，算出传质因子 $j_D$，据此判断在该转速（或循环量）下颗粒外表面与流体之间是否
还有可观的浓度差。（2）\textbf{返混实验}：在一定流速的进器气流中注入示踪气脉冲，记录
出器气流中示踪物浓度随时间的应答曲线，由曲线形状测出返混程度，确认「器内组成等于出口
组成」这一 CSTR 基本假定成立。两项都通过，转速（循环量）才算选对。
\src{S2 §10.5.1，书内 p.783—784 / PDF p.272—273}
\end{measure}

\begin{pitfall}
\textbf{同一种催化剂、同一个反应，在搅拌（循环）反应器与管式反应器中可能给出不一致的
动力学结果}——这未必是实验做错了，而可能是催化剂处于\textbf{不同的活性状态}。原书给出
的实例：邻二甲苯在两种反应器中的结果显示，搅拌反应器中催化剂的再氧化过程起控速作用，
催化剂处于还原态；而在进料组成相同的管式积分反应器中，烃的外扩散起控速作用，催化剂
处于氧化态，因而活性不同。\textbf{推论}：把非稳态信息（表面状态）与稳态动力学（速率
方程）分开报告，一旦两者矛盾，先查控速步骤是否在两个装置中换了人。
\src{S2 §10.5.1，书内 p.784 / PDF p.273}
\end{pitfall}

\begin{crosscheck}
选型表只回答「能不能测」，不回答「测得准不准」。表的第四列里凡是出现「需检验」「须
假定」的条目（无梯度反应器的传质与返混检验、绝热床的梯度可忽略假定、微分床的分析精度），
都必须在实验记录中留下检验证据；否则该行应从「能测」降级为「存疑」。
与 10.4 节的传递过程判据（Da 数、有效因子）配合使用，可把这里的定性判断变成定量判断。
\end{crosscheck}

\subsubsection{内扩散参数的测定}

扩散系数 $D_B$、$D_K$、$D_e$ 既有估算法，也有无反应下的实验测定法。测定方法上过渡应答
类技术发展较快，如产物瞬态分析（TAP）技术，应答时间可达毫秒量级；由脉冲技术的动力学
数据可用矩量分析、Fourier 拟合、Laplace 拟合等多种方法拟合扩散系数与平衡常数。
但必须注意：\textbf{无反应下估算或测定的数值与有反应情况下的实际有效扩散系数差距较大，
甚至可达 4—5 倍}——原因是反应与非反应两种情况下扩散通量比相差甚大，此外还与表面扩散
作用、孔径分布不匀有关。因此\textbf{在反应情况下测得的数据求内扩散参数更妥当}。
\src{S2 §10.5.2，书内 p.784—785 / PDF p.273—274}

\begin{principle}
\textbf{方法一：有效因子法。}在恒定条件下测定不同粒径催化剂的反应速率。当粒径小到
反应速率不再变化时，该极限速率对应的 $\eta=1.0$；某一粒径催化剂上的速率与极限速率之比
即为该粒径下的有效因子 $\eta$。对多种粒径取样测定，可得一组 $\eta$ 值，进而求出扩散
系数。若只有两种粒径（等效直径 $d_1$、$d_2$，速率 $r_1$、$r_2$），可用所谓的
\textbf{「三角形」方法}：利用 $\eta$ 与内扩散模数 $\phi$ 的关系在一定范围内比较固定的
特点，有
\begin{equation}
\frac{r_1}{r_2}=\frac{\eta_1}{\eta_2},\qquad
\frac{d_1}{d_2}=\frac{\phi_1}{\phi_2}
\label{eq:tri-method}
\end{equation}
在 $\eta$–$\phi$ 图上，$\eta_1/\eta_2$ 是纵坐标上的一段长度、$\phi_1/\phi_2$ 是横坐标上
的一段长度，二者分别是某一直角三角形的两条直角边；保持直角边与坐标轴平行，把三角形
顶点沿 $\eta$–$\phi$ 曲线滑移，直到两个顶点同时落在曲线上，即读出相应的 $\eta_1$、
$\eta_2$、$\phi_1$、$\phi_2$。注意\textbf{待拟合的这段曲线不在 $\eta$–$\phi$ 的直线段上}。
\src{S2 §10.5.2，书内 p.785 / PDF p.274}
\end{principle}

\begin{example}[两粒径三角形法求有效因子（丁烯氧化脱氢）]
同一催化剂取两种颗粒半径 $0.025$~cm 与 $0.160$~cm，测得丁烯转化速率：$410$\,\degC 下
分别为 $8.89$ 与 $8.04$~mmol\,g$^{-1}$\,h$^{-1}$；$470$\,\degC 下分别为 $14.25$ 与
$10.89$~mmol\,g$^{-1}$\,h$^{-1}$。于是 $\phi_1/\phi_2=d_1/d_2=0.160/0.025=6.4$，
而 $\eta_2/\eta_1=r_2/r_1$ 为 $1.106$（$410$\,\degC）与 $1.308$（$470$\,\degC）。
按式~\eqref{eq:tri-method} 在 $\eta$–$\phi$ 图上作横坐标长 $6.4$、纵坐标长 $1.106$
（或 $1.308$）的直角三角形，滑移后读出：$410$\,\degC 时 $\phi_1=1.30$、$\eta_1=0.905$，
$\phi_2=0.20$、$\eta_2=1.00$；$470$\,\degC 时 $\phi_1=2.30$、$\eta_1=0.757$，
$\phi_2=0.36$、$\eta_2=0.99$。有了 $\eta$ 值，体系受扩散控制的程度就有了量的概念，
并可进一步求有效扩散系数 $D_e$。\src{S2 §10.5.2，书内 p.785—786 / PDF p.274—275}
\end{example}

\begin{principle}
\textbf{方法二：Weisz 内扩散模数 $\Phi$。}为便于求 $D_e$，Weisz 定义了一个新的内扩散
模数
\begin{equation}
\Phi=\eta\phi^{2}=\frac{d^{2}}{D_e}
\left(\frac{1}{V_p}\frac{\dd n}{\dd t}\right)\frac{1}{c_s}
\label{eq:weisz-modulus}
\end{equation}
式中球形模型取 $d=R$，平板模型取 $d=L$。$\eta$–$\Phi$ 关系曲线不像 $\phi$ 那样包含
化学反应本身的速率常数与反应级数，因此可以从实测 $\eta$ 先读出 $\Phi$，再按
式~\eqref{eq:weisz-modulus} 求 $D_e$，比用 $\phi$ 方便。
\tocheck{未对照原页面图像：式~\eqref{eq:weisz-modulus} 中体积项与速率项的归并方式
（$1/V_p$ 是否置于括号内）以及消耗速率是否取负号，OCR 无法判定；数值关系
$\Phi=\eta\phi^{2}$ 本身可确认。}\src{S2 §10.5.2，书内 p.785—786 / PDF p.274—275}
\end{principle}

\textbf{方法三：单锭扩散反应器（single pellet diffusion reactor）。}把催化剂锭片挤压在
不锈钢柱体中（压锭时使密度分布尽可能均匀），反应气体以足够速度用往复泵压经锭片正面，
用注射针筒取样作色谱分析；一部分反应气体扩散到锭片背面，也取样分析（也可用质谱或
红外）。取样量很少、催化剂活性稳定、到达定态后，按 Petersen 的理论分析，对不同反应
级 $m$，锭片背面浓度 $c_1(1)$ 与正面浓度 $c_1(0)$ 之比 $\Psi_1$ 与 $c_1(0)$ 之间有
确定的函数关系，由此可由一系列 $(\Psi_1,\,c_1(0))$ 实验值定出反应级数 $m$；已知 $m$ 后，
再由一次实验的 $\Psi_1$ 与 $\phi$ 的函数关系求出内扩散模数 $\phi$，进而得 $D_e$。
原书对 $m=2,\,1/2,\,0,\,-1/2,\,-3/4,\,-1,\,-2$ 等情形给出了相应曲线。
\src{S2 §10.5.2，书内 p.786—787 / PDF p.275—276}
\tocheck{未对照原页面图像：$\Psi_1$ 与 $c_1(0)$、$\Psi_1$ 与 $\phi$ 的两组关系曲线
（原书图 10-14、图 10-15）为读图求值用，本分片无法核对曲线坐标与刻度。}

\begin{dataflow}
\textbf{从「内扩散参数」到「能用哪一档速率数据」的判断链}：（1）同一批催化剂取 3—5 种
粒径，在\textbf{同一空速、同一温度}下测转化率；（2）转化率随粒径减小而上升并趋于极限值
$\Rightarrow$ 小粒径区已排除内扩散，取该粒径以下的数据作动力学；（3）转化率始终随粒径
变化而未见平台 $\Rightarrow$ 尚未脱离内扩散控制，加做更小粒径或改换孔结构；（4）把极限
粒径下的 $\eta$ 记为 1.0，按方法一求各粒径的 $\eta$；（5）用式~\eqref{eq:weisz-modulus}
得 $\Phi$，与判据阈值（见 10.5.4 的式~\eqref{eq:wp}）比较，决定数据是「本征」还是
「表观」。\src{S2 §10.5.2，书内 p.785—786 / PDF p.274—275}
\end{dataflow}

\subsubsection{速率方程中吸附参数的测定：过渡应答法}

速率方程里的吸附平衡常数 $K_i$ 是按 Langmuir 图景推出的，若不能在\textbf{反应条件下}
测定吸附，机理与速率控制步骤的讨论就始终悬空。\textbf{过渡应答法}
（transient response method）正是为此发展起来的：对一个稳态反应系统施加某种微扰，
系统在趋近新稳态过程中的\textbf{应答方式对不同机理各有特征}，由此可以抽提出吸附物种的
吸附量、吸附与脱附速率的信息。对多相催化，施加\textbf{浓度微扰}最为方便。
\src{S2 §10.5.3，书内 p.788 / PDF p.277}

\begin{principle}
以一个双组分反应 $\mathrm{B_1+B_2\to B_3}$ 为例，设它由三个基元步骤构成：
$\mathrm{B_1+L\rightleftharpoons B_1L}$（$k_1$、$k_{-1}$）、
$\mathrm{B_2+B_1L\rightleftharpoons B_3L}$（$k_2$、$k_{-2}$）、
$\mathrm{B_3L\rightleftharpoons B_3+L}$（$k_3$、$k_{-3}$）。在理想流体（全返混或活塞流）
且无传热、传质限制的前提下，当 $\mathrm{B_1}$、$\mathrm{B_2}$ 与大量惰性气以恒定组成、
恒定流速通过床层达稳态后，令 $\mathrm{B_1}$（或 $\mathrm{B_2}$、或 $\mathrm{B_3}$）的
浓度发生一个明确的\textbf{阶跃}变化（分压自 $p_{10}$ 跃至 $p_{1s}$）。
\textbf{（1）吸附量与吸附速率的读出}：若进料中原先无 $\mathrm{B_2}$，$\mathrm{B_1}$ 的
分压沿一条曲线渐升至 $p_{1s}$，此时曲线与终值之间的\textbf{阴影面积} $Q_1$ 即自 $t=0$ 至
$t_1$ 时间内 $\mathrm{B_1}$ 所增加的吸附量，而 $t_1$ 时刻 $\mathrm{B_1}$ 的吸附速率由
$\Delta p_1F/p$ 给出（$F$ 为进料速率，$p$ 为总压）；若进料中原已有 $\mathrm{B_2}$，
$\mathrm{B_1}$ 沿另一条曲线渐升至较低的 $p'_{1s}$（$p'_{1s}<p_{1s}$）。
\textbf{（2）速率控制步骤的判据}：引入 $\mathrm{B_1}$ 阶跃后，若产物 $\mathrm{B_3}$ 的
分压应答显示\textbf{立即}到达新的稳态，说明 $\mathrm{B_1}$ 的吸附或脱附在动力学上都不
重要，速率控制步骤是表面反应。\textbf{（3）可逆性的判据}：把进料中 $\mathrm{B_1}$ 与
$\mathrm{B_2}$ 的分压都突降至零，观察 $\mathrm{B_3}$ 的应答——$k_{-2}\approx 0$（步骤 2
不可逆）时曲线偏高，曲线下的阴影面积 $Q_2$ 是到 $t_1$ 时 $\mathrm{B_3}$ 的吸附量、
$\Delta p_3F/p$ 是其净脱附速率；步骤 2 可逆时曲线明显偏低。
\src{S2 §10.5.3，书内 p.788—789 / PDF p.277—278}
\end{principle}

\begin{measure}
\textbf{过渡应答实验的落点。}物料衡算给出一个微分方程组，其解给出一族以
$k_1$、$k_2$、$k_3$、$k_{-3}$ 为参数的应答曲线。原书给出一组具体条件：长为 $100$~cm、
催化剂堆密度 $1.0$~g\,cm$^{-3}$、空隙率 $0.5$ 的微分反应管，以线速 $400$~cm\,min$^{-1}$
通入 $20.3$~kPa 的 $\mathrm{B_1}$ 与 $80.1$~kPa 氦的混合气达稳态，然后施加浓度
$20.3$~kPa 的 $\mathrm{B_2}$ 阶跃，记录 $\mathrm{B_3}$ 的过渡应答曲线（取
$k_{-1}=k_{-2}=0$）。曲线形状随参数组合系统变化：当 $k_1\approx k_2\gg k_3$ 时得到一类
典型结果；当 $k_1$ 较高而 $k_2$ 逐步接近 $k_3$ 时得到另一类；而当 $k_1$ 很小，
$p_3$ 会很快升到一高值后旋即随时间渐降——因为经基元步骤 2 所消耗的吸附 $\mathrm{B_1}$
来不及补充。\textbf{操作要点}：反应器的自由体积、取样管线与检测器响应时间共同构成一个
额外的「仪器应答函数」，必须先用空白阶跃测出它，否则这族曲线无法与理论曲线对照。
\src{S2 §10.5.3，书内 p.789 / PDF p.278}
\end{measure}

过渡应答法用以研究 CO 催化氧化和 NO 催化分解等简单反应取得了巨大的成功；在丙烯氧化、
F–T 合成、甲烷氧化等反应中也得到很好结果。随着脉冲进样技术与检测技术的提高，它正成为
一种重要的动力学研究技术。\src{S2 §10.5.3，书内 p.789 / PDF p.278}

\begin{crosscheck}
过渡应答法给出的是\textbf{吸附量、吸附／脱附速率与控速步骤的直接证据}，而稳态动力学
给出的是\textbf{速率对流体相组成的依赖}——两者不能互相替代：稳态数据只能定出 $K_i$ 的
组合（例如双曲线式分母中 $K_1c_1+K_2c_2$ 的各项），无法区分「$K$ 大而覆盖度高」与
「$K$ 小而覆盖度低」；过渡应答则直接测吸附量。反过来，过渡应答的解析依赖一个假定的
基元步骤网络，网络错了，拟合出的速率常数仍可能看起来很好——所以它必须与稳态的级数、
活化能互相回查。更短时间尺度上的同类信息由 SSITKA 与 TAP 给出（见 10.6.3、10.6.4）。
\end{crosscheck}

\subsubsection{本征动力学区的确定与传递过程干扰的排除}

对多相催化反应，随温度升高可以观察到三个不同的催化反应区域。
\textbf{（1）本征动力学区}：温度足够低时反应速率很慢，扩散所需的势能可以忽略，此时
观察到的是反应本征动力学。\textbf{（2）内扩散控制区}：温度升高后扩散速率增加得较慢，
而本征速率常数呈指数增长，扩散消耗掉更多的有效势能，只有少量用以驱动化学反应，颗粒内
产生显著的反应物浓度梯度。\textbf{（3）外扩散控制区}：温度进一步升高，有效因子 $\eta$
逐渐下降，反应物在进入颗粒内部之前就被消耗尽，流体体相与颗粒外表面之间的浓度差变大
并逐渐重要。三个区域之间的相对位置与众多因素有关；在内扩散区，速率可借减小颗粒尺寸与
改变孔结构来提高，\textbf{但表观活化能不变}；使用足够小的颗粒且为无孔催化剂时，内扩散
可以消失，出现本征动力学区向外扩散区的\textbf{直接过渡}。\textbf{在外扩散区，反应速率
与流体线速度有直接的函数关系；在内扩散区，反应速率与线速度无关}——这两句话是所有
排除实验的出发点。\src{S2 §10.5.4，书内 p.790—791 / PDF p.279—280}

\begin{keycond}
\textbf{为什么必须排除外扩散与内扩散。}（1）速率方程是要外推到工业粒度、工业线速和
工业温度去的，若式中的「速率」还含设备参数（线速、粒径），外推在原理上就不成立。
（2）表观量会被系统性地歪曲：强内扩散控制下一级反应的表观活化能约为真实值的一半
（$E_{\mathrm{app}}\approx E/2$），表观级数也被推离真值，用这样的 $E$ 外推温度效应
必然出错。（3）判据的意义在于区分「没排除」与「已排除」是可检验的：前者表现为速率随
线速或粒径变化，后者表现为速率对二者不敏感。传递过程对活化能与级数影响的具体推导见
10.4.2 节。\src{S2 §10.4.2，书内 p.772—775 / PDF p.261—264；§10.5.4，书内 p.790—792 /
PDF p.279—281}
\end{keycond}

\textbf{要把「排除扩散影响」变成可操作的实验，只有三组自变量}：\textbf{线速（或质量
流速）、粒径、温度}。下面按原书给出的做法逐条列出，并给出定量判据。

\textbf{实验一：在相同停留时间下改变流速（外扩散）。}在同一反应器内，随反应物的质量
流速变化而相应增减催化剂用量，\textbf{以保持相同的接触时间}；若质量流速增大而转化率趋于
稳定值，则外扩散已消除。另一种做法基于相同思想：采用截面积相同的两个反应器，填充不同量
的催化剂，引入与催化剂量成正比的反应物流量，使两者接触时间相同；高流速区两反应器转化率
无差别，低流速区则不同。一般应在流速高到它对转化率无明显影响的区域内取动力学数据。
还可以在同一个反应器中改变流量、在不同温度下实验并计算表观速率常数：若不同流量的曲线
\textbf{重合}，则外扩散不存在；若在较高温度下几条曲线\textbf{不重合}，表示相应于不同
流量存在不同程度的外扩散。\src{S2 §10.5.4，书内 p.791—792 / PDF p.280—281}
\textbf{注意用词}：此处的「相同空速」指接触时间 $W/F$ 不变，\textbf{不是}指体积流量不变；
把流量与催化剂量同比改变，才是在检验线速的影响。

\textbf{实验二：改变催化剂粒径（内扩散）。}从同一批催化剂取得不同粒径的几份样品，
测同一空速下的转化率；一般随粒径减小，转化率上升到一极限值。取与该极限值相应的粒径
以下（更小）的催化剂作研究，一般认为即可排除内扩散干扰。
\src{S2 §10.5.4，书内 p.792 / PDF p.281；图 10-22，书内 p.793 / PDF p.282}

\textbf{实验三：变更温度重复实验一、二（把三区画出来）。}只有在前两组实验于\textbf{整个
实验温度区间}都通过，才可以在该区间内称数据为本征动力学数据；先在高温区通过、
在低温区没有检验过，同样不可外推。

\begin{principle}
\textbf{三个定量判据。}\textbf{前提：三式里的同一个符号 $r$ 必须指同一种速率。}
动力学数据常以\textbf{单位质量}报出，而三个判据本来是按\textbf{单位颗粒体积}推出来的，
两者差一个颗粒密度：
\begin{equation}
r_v=\rho\,r_m,\qquad
r_m\ \text{为}\ \unit{mol.kg^{-1}.s^{-1}}\ \text{（按催化剂质量）},\quad
r_v\ \text{为}\ \unit{mol.m^{-3}.s^{-1}}\ \text{（按颗粒体积）}
\label{eq:rate-basis}
\end{equation}
式中 $\rho$ 为\textbf{颗粒密度}（含内孔，不是骨架密度）。下文三式中的 $r$ 一律指
$r_v$，即\textbf{每立方米颗粒每秒转化的摩尔数}；手中只有质量速率时要按
$r\to\rho r_m$ 替换后再代入（三式的质量速率写法逐一列在式后）。
此外三个判据的 $R$ 都是\textbf{颗粒半径}，
但传热判据里的 $R T_B/E$ 要求 $R$ 是\textbf{气体常量}——同一字母两种含义，
代入时按单位自行区分（见下方误判提示）。$R$ 与几何的关系：对球形颗粒，由
$S_{\text{外}}/V_p=3/R$ 反解得 $R=3V_p/S_{\text{外}}$，即\textbf{半径取体积与外表面积之比
的三倍}；非球形颗粒用同一关系折算成等效半径，不要直接拿筛分粒度。

\textbf{（1）外扩散（Damköhler 数判据）}：对 $m$ 级反应，
\begin{equation}
Da=\frac{rR}{k_m c_B}<\frac{0.15}{m}
\label{eq:da}
\end{equation}
式中 $c_B$ 为组分 $\mathrm{B_1}$ 在流体相的浓度，$r=r_v$ 为按颗粒体积计的观测速率，
$R$ 为催化剂颗粒半径，$k_m$ 为颗粒外膜的传质系数。
\textbf{量纲}：$\unit{mol.m^{-3}.s^{-1}}\times\unit{m}\big/
\big(\unit{m.s^{-1}}\times\unit{mol.m^{-3}}\big)$——无量纲。
改用质量速率时本式写成 $Da=\rho r_mR/(k_mc_B)$，与上式\textbf{数值恒等}。
\textbf{（2）内扩散（Weisz–Prater 判据）}：
% 【量纲修正 2026-09-29】原式写作 \Phi = r\rho R^2/(D_e c_B)，并声明「r 一律指 r_v」。
% 但 r_v 已是体积基准 mol.m^-3.s^-1，再乘 rho(kg.m^-3) 后整体量纲是 **kg.m^-3**，
% 不是无量纲——本节自己的量纲式也得出 kg.m^-3，却写成了「无量纲」。
% 随后又说质量基准要写成 r_m\rho^2 R^2/(...)——那是把 rho 记了两次。
% 正确形式（与第 10 章分片 b 及评审 C10-02 一致）：
%   体积基准 Phi = r_v R^2/(D_e c_B)        —— 无 rho
%   质量基准 Phi = rho r_m R^2/(D_e c_B)    —— 恰好一个 rho
% 原书印成「r rho R^2/(...)」只有在 r 取**质量基准**时才成立
% （原书数据本身就是 mmol/(g.h)）——这就是评审指出的缺陷根源。
\begin{equation}
\Phi=\frac{r_v R^{2}}{D_e c_B}
\label{eq:wp}
\end{equation}
式中 $r_v$ 为\textbf{体积基准}的反应速率，$D_e$ 为有效扩散系数。\textbf{量纲}：
$\unit{mol.m^{-3}.s^{-1}}\times\unit{m^2}\big/
\big(\unit{m^2.s^{-1}}\times\unit{mol.m^{-3}}\big)$——无量纲。
改用\textbf{质量基准}速率 $r_m$ 时，用 $r_v=\rho r_m$ 代入，本式写成
$\Phi=\rho r_mR^{2}/(D_ec_B)$，与上式\textbf{数值恒等}：$\rho$ 在质量基准形式里
\textbf{只出现一次}。
\tocheck{原书式（10-131）印为「$r\rho R^{2}/(D_ec_B)$」而无量纲式又给出
$\unit{kg.m^{-3}}$；按上述两式对照，原书的 $r$ 应取\textbf{质量基准}
（其数据本身即 $\unit{mmol.g^{-1}.h^{-1}}$），或原式多写了一个 $\rho$。
本节按体积基准形式给出，与式~\eqref{eq:da} 及第 10 章分片 b 的判据统一；
两种基准相差 $\rho_p$ 约三个数量级，用错会让临界粒径差上千倍。}
对等温球形颗粒，一级反应时 $\Phi<1$ 对应有效因子约 $0.94$；二级反应取 $\Phi<0.3$。
\textbf{（3）相间传热（Mears 判据）}：$\eta>0.95$ 的判据是
\begin{equation}
\frac{r(-\Delta H)R}{h_BT_B}<0.15\,\frac{\mathscr{R}T_B}{E}
\label{eq:mears-heat}
\end{equation}
式中 $T_B$ 为流体温度，$h_B$ 为相间传热系数，$E$ 为活化能；为消除同一字母的歧义，
\textbf{右侧的气体常量记作 $\mathscr{R}$}，左侧的 $R$ 仍是颗粒半径。
\textbf{量纲}：$\unit{mol.m^{-3}.s^{-1}}\times\unit{J.mol^{-1}}\times\unit{m}\big/
\big(\unit{W.m^{-2}.K^{-1}}\times\unit{K}\big)$——无量纲；右侧
$\unit{J.mol^{-1}.K^{-1}}\times\unit{K}\big/\big(\unit{J.mol^{-1}}\big)$ 亦无量纲。
改用质量速率时左侧写成 $\rho r_m(-\Delta H)R/(h_BT_B)$，同样\textbf{数值恒等}。
\src{S2 §10.5.4，书内 p.791—792 / PDF p.280—281；§10.5.4 末段，书内 p.793 / PDF p.282}
\end{principle}

\begin{pitfall}
\textbf{三个判据必须用同一速率基准，漏乘或重乘 $\rho$ 都会让判据失效。}判据 (1) 与 (3)
的分子必须是\textbf{体积速率} $r_v$（$\unit{mol.m^{-3}.s^{-1}}$）；把催化研究最常用的
$\unit{mol.g^{-1}.s^{-1}}=\unit{mol.kg^{-1}.s^{-1}}\times10^{3}$ 直接当 $r$ 代进去，
得到的"判据"单位是 $\unit{m^3.kg^{-1}}$——\textbf{有量纲、无论数值大小都不成判据}，
这正是"看起来通过判据"的假阴性来源。自检办法只有一个：
\textbf{逐式展开 SI 单位，确认结果是纯数}（如上三式）。
注意判据 (2) 与 (1)(3) 形式上不同，不是因为它换了基准，而是因为它的分子里本来就有
一个 $\rho$（把表面速率换算成体积基准时引入）：把 $r_m$ 代进 (2) 时要写成
$r_m\rho^{2}R^{2}/(D_ec_B)$ 才与 $r_v\rho R^{2}/(D_ec_B)$ 相等，
若只写成 $r_m\rho R^{2}/(D_ec_B)$ 就少了 $\rho$ 倍、把小颗粒误判成受内扩散控制。

\textbf{式~\eqref{eq:wp} 两处印刷／OCR 存疑，务必按原页面核验后再引用。}
（a）原书表述为「一级反应时 $\Phi<1$ 可保证内扩散有效因子小于 $0.95$」，但按球形颗粒
一级反应的精确关系 $\eta=(3/\phi)(1/\tanh\phi-1/\phi)$ 复核，$\Phi=\eta\phi^{2}=1$ 给出
$\eta\approx 0.94$，即 $\Phi<1$ 对应 $\eta$ \textbf{大于}约 $0.94$；原文的「小于」应为
「大于」或为 OCR 失真。（b）末句「为保险起见，内扩散影响排除一般要求 $\Phi$ ? 1」中的
不等号在 OCR 中丢失，符号与阈值都无法确认。
\tocheck{未对照原页面图像：式~\eqref{eq:wp} 的阈值方向（$\eta$ 大于还是小于 $0.95$）
与末句「$\Phi$ 取何值才算保险」的数值，均需查原书 p.792。}
\textbf{式~\eqref{eq:mears-heat} 的符号冲突}：不等式左侧的 $R$ 是颗粒半径，右侧的 $R$
是气体常数，原书同一式中两者同形；本稿把右侧改写为 $\mathscr{R}$ 以消除歧义，原书写法
仍是两个 $R$。\tocheck{未对照原页面图像：式~\eqref{eq:mears-heat}
右侧 $RT_B/E$ 中的 $R$ 是否为气体常数、$h_B$ 是否为相间传热系数，需核验；量纲上
只有当右侧 $R$ 为气体常数时两侧才无量纲。}
\src{S2 §10.5.4，书内 p.792—793 / PDF p.281—282}
\end{pitfall}

\textbf{同一组数据换基准，判据必须不变（数值自检）。}取颗粒密度
$\rho=\qty{1500}{kg.m^{-3}}$、$R=\qty{1.0}{mm}$、$k_m=\qty{0.10}{m.s^{-1}}$、
$c_B=\qty{10}{mol.m^{-3}}$、$D_e=\qty{1.0e-6}{m^2.s^{-1}}$、
$\Delta H=\qty{-1.0e5}{J.mol^{-1}}$、$h_B=\qty{100}{W.m^{-2}.K^{-1}}$、
$T_B=\qty{500}{K}$、$E=\qty{1.0e5}{J.mol^{-1}}$，并设实测
$r_m=\qty{1.0e-4}{mol.kg^{-1}.s^{-1}}$，则 $r_v=\rho r_m=\qty{0.15}{mol.m^{-3}.s^{-1}}$。
三条判据分别算两次——一次用 $r_m$ 配密度、一次直接用 $r_v$：

（i）\textbf{Damköhler}：$r_m\rho R/(k_mc_B)
=10^{-4}\times1500\times10^{-3}/(0.10\times10)=1.5\times10^{-4}$，与
$r_vR/(k_mc_B)=0.15\times10^{-3}/(0.10\times10)=1.5\times10^{-4}$ 完全相同，
两者都\textbf{无量纲}。

（ii）\textbf{Weisz--Prater}：以体积速率代入
$r_v\rho R^{2}/(D_ec_B)=0.15\times1500\times10^{-6}/(10^{-6}\times10)
=2.25\times10^{1}$；以质量速率代入
$r_m\rho^{2}R^{2}/(D_ec_B)=10^{-4}\times1500^{2}\times10^{-6}/(10^{-5})
=2.25\times10^{1}$。两者仍是同一个数——\textbf{但前提是分子上有两个 $\rho$}。
若只写一个 $\rho$（$r_m\rho R^{2}/(D_ec_B)=1.5\times10^{-2}$），
结果就小了 $1500$ 倍，$\Phi=22.5$ 会被误算成 $0.015$、从「受内扩散控制」变成
「已排除内扩散」，结论完全相反。

（iii）\textbf{Mears}：$r_m\rho(-\Delta H)R/(h_BT_B)
=10^{-4}\times1500\times10^{5}\times10^{-3}/(100\times500)=3.0\times10^{-4}$，
与 $r_v(-\Delta H)R/(h_BT_B)=0.15\times10^{5}\times10^{-3}/(5\times10^{4})
=3.0\times10^{-4}$ 相同，两者都\textbf{无量纲}。

反过来，若把 $\qty{1.0e-4}{mol.kg^{-1}.s^{-1}}$ \textbf{不乘密度}直接当 $r$ 代进
$Da$ 与 Mears 式，得到的是 $\qty{1e-7}{m^3.kg^{-1}}$、$\qty{2e-7}{m^3.kg^{-1}}$
这样的\textbf{有量纲}量，无论大小都不能与阈值比较——这正是「看起来通过判据」的假阴性来源。

\textbf{活塞流前提的几何条件。}在保证理想混合模型与理想活塞流状态方面，原书给出的条件
是：反应管内径应至少为催化颗粒直径的 $10$ 倍以上；对气-固反应，催化剂床层长度至少应为
催化剂粒径的 $50$ 倍以上；对液-固反应，床层长度至少应为催化剂粒径的几百倍以上。
管径过小而颗粒很大、特别是装填不匀时，会发生沟流，使垂直于流向的截面内各部分表观线速
不均一；床层过短则轴向弥散严重。两者都会破坏活塞流假定。
\src{S2 §10.5.4，书内 p.791 / PDF p.280}

\begin{figure}[htbp]
\centering
\begin{tikzpicture}
\tikzset{declare function={
  xint(\T)=0.92*(44000000*exp(-9622/\T))/(1+44000000*exp(-9622/\T));
  xext(\u)=2.5*sqrt(\u/30);
  xobs(\u,\T)=1/(1/xint(\T)+1/xext(\u));
}}
\begin{axis}[
  width=\onecolwidth, height=52mm,
  xlabel={线速 $u$ / (cm$\cdot$min$^{-1}$)},
  ylabel={转化率 $x$},
  domain=10:800, samples=140,
  xmin=10, xmax=800, ymin=0, ymax=0.95,
  legend pos=south east,
  legend cell align=left,
  grid=major,
  tick label style={font={\fontsize{7.5}{9}\selectfont}},
  label style={font={\fontsize{8}{9.5}\selectfont}},
  legend style={font={\fontsize{7}{9}\selectfont}, fill=none},
]
\addplot[very thick] {xobs(x,470)};
\addlegendentry{$T=470$~K：曲线水平，外扩散不存在}
\addplot[very thick, dashed] {xobs(x,520)};
\addlegendentry{$T=520$~K：仍有可辨的线速依赖}
\addplot[very thick, dotted] {xobs(x,570)};
\addlegendentry{$T=570$~K：外扩散起作用，趋向平台}
\end{axis}
\end{tikzpicture}
\caption{模型计算：\textbf{接触时间（空速）保持不变}、只改变线速时，转化率随线速的变化
如何判别外扩散控制。读者应看三条曲线的\textbf{形状差异}：低温曲线在整个线速区间水平，
说明速率与线速无关，数据落在本征动力学区；温度升高后曲线左端明显偏低、随线速上升并
逐渐逼近平台，左端即外扩散控制区，右端的平台才是可以取数的区域。\textbf{条件与局限}：
等温、无内扩散（取足够小的粒径）、反应不可逆且远离平衡；模型为串联阻力式
$1/x=1/x_{\text{本征}}+1/x_{\text{外}}$，其中外扩散可达的极限转化率取
$x_{\text{外}}\propto k_m\propto u^{0.5}$、本征可达转化率取 Arrhenius 型有界式
$x_{\text{本征}}=0.92\,A/(1+A)$、$A=4.4\times10^{7}\exp(-9622/T)$；纵轴数值为示意，
不代表任何具体反应体系；三条曲线的温度为 $470$、$520$、$570$~K，绝对位置可整体缩放而
不改变判别结论。\modelmark\figsrc{据 S2 §10.5.4 图 10-20、图 10-21 的判据改绘为模型计算}}
\label{fig:ch10-velocity}
\end{figure}

\begin{figure}[htbp]
\centering
\begin{tikzpicture}
\tikzset{declare function={
  etaWP(\p)=(3/\p)*(1/tanh(\p)-1/\p);
  PhiWP(\p)=etaWP(\p)*\p*\p;
}}
\begin{axis}[
  width=\onecolwidth, height=50mm,
  xmode=log,
  xlabel={Weisz 模数 $\Phi=\eta\phi^{2}$},
  ylabel={有效因子 $\eta$},
  domain=0.3:20, samples=160,
  ymin=0, ymax=1.08,
  legend pos=south west,
  grid=major,
  tick label style={font={\fontsize{7.5}{9}\selectfont}},
  label style={font={\fontsize{8}{9.5}\selectfont}},
  legend style={font={\fontsize{7}{9}\selectfont}, fill=none},
]
\addplot[very thick] ({PhiWP(x)},{etaWP(x)});
\addlegendentry{球形颗粒、一级反应（精确解）}
\addplot[dashed, domain=0.3:60] {0.95};
\addlegendentry{$\eta=0.95$ 参考线}
\end{axis}
\end{tikzpicture}
\caption{模型计算：Weisz 模数与有效因子的关系，用于把式~\eqref{eq:wp} 的判据落到可读的
数值上。读者应看\textbf{曲线与 $\eta=0.95$ 参考线的交点}：交点位于 $\Phi\approx 0.75$
—$1.0$ 之间，即判据取 $\Phi<1$ 时有效因子不低于约 $0.94$；$\Phi$ 再大一个数量级，
$\eta$ 已降到 $0.5$ 以下，此时测得的活化能与级数都不能当本征值用。\textbf{条件与局限}：
曲线由球形颗粒一级反应的精确解 $\eta=(3/\phi)(1/\tanh\phi-1/\phi)$ 与
$\Phi=\eta\phi^{2}$ 数值算出，只适用于等温、一级、球形；其它级数与形状的曲线位置不同，
故原书对二级反应另给 $\Phi<0.3$。\modelmark\figsrc{据 S2 §10.5.2 图 10-12 与 §10.5.4
式（10-131）重算}}
\label{fig:ch10-etaphi}
\end{figure}

\begin{dataflow}
\textbf{「本征动力学数据」的取得流程与失败分支。}
（1）按几何条件设计反应管：内径 $\ge 10d_p$，气固床层长 $\ge 50d_p$（液固为几百倍
$d_p$）。（2）用萘锭传质实验与示踪脉冲返混实验确认装置的无梯度性（无梯度反应器）
或活塞流性（管式反应器）。（3）做实验一（变线速、保持 $W/F$）：转化率出现平台
$\Rightarrow$ 外扩散已排除；用式~\eqref{eq:da} 复核 $Da<0.15/m$。（4）做实验二
（变粒径、保持空速）：转化率出现极限值 $\Rightarrow$ 取极限粒径以下的数据；用
式~\eqref{eq:wp} 复核 $\Phi<1$（二级反应 $<0.3$）。（5）用
式~\eqref{eq:mears-heat} 确认相间温差可忽略。（6）在通过（3）—（5）的温度区间内拟合
速率方程；任一步失败，就只能报告\textbf{表观}速率方程，并注明受控的传递环节。
\src{S2 §10.5.4，书内 p.790—793 / PDF p.279—282}
\end{dataflow}

\begin{pitfall}
\textbf{本节的四个高频误判。}
（1）\textbf{在传质控制区测「本征活化能」，得到偏低的假值。}强内扩散控制下一级反应的
表观活化能约为真实值的一半（$E_{\mathrm{app}}\approx E/2$），表观级数同样被推离真值。
判别方法：看 $E$ 的数值是否随粒径变化——改变粒径而 $E$ 变，说明测的不是本征值；
在内扩散区减小粒径可提高速率，但\textbf{不会}改变真实活化能这一事实，恰是分解二者的
钥匙。\src{S2 §10.4.2，书内 p.772—775 / PDF p.261—264}
（2）\textbf{催化剂装填量过大导致床层温度梯度。}积分反应器中转化率高、反应热稍大就
难恒温；装填量越大、床层越长，轴向温差越大，而速率对温度是指数敏感的，于是「床层平均
温度」变成一个人造量。对策：减小管径、恒温导热介质、惰性稀释、非等比稀释
（入口加大稀释比），必要时改用微分床或无梯度反应器。\src{S2 §10.5.1，书内 p.779—780 /
PDF p.268—269}
（3）\textbf{转化率取点太靠近平衡。}可逆反应的速率中含推动力项 $(c_1-c_2/K_{eq})$，
当 $c_2/c_1\to K_{eq}$ 时推动力趋于零，表观反应级数趋于零、速率趋于零，此时测得的
「级数」与「活化能」都只是分析误差的放大器，与机理无关。取点应保证
$c_2/c_1$ 明显小于 $K_{eq}$，并在报告中给出该比值。
（4）\textbf{产物分析滞后导致非稳态数据错配。}SSITKA、TAP、过渡应答测的都是\textbf{瞬变
曲线}，而色谱一轮分析以分钟计，用它跟瞬变过程对不上：进样阀、管线和检测器死体积会把
阶跃展宽成一个「仪器应答函数」。正确做法是用多通道快速质谱（这正是 SSITKA 与 TAP 都对
检测器提出「时间分辨率高」「电离室距反应器出口 $\le 1$~cm」要求的理由），并先用空白
阶跃测出仪器应答，再把实测曲线与理论曲线对照；色谱只用于测定稳态浓度。
\src{S2 §10.6.3，书内 p.798 / PDF p.287；§10.6.4，书内 p.799 / PDF p.288}
\end{pitfall}

\subsection{非稳态催化过程动力学}

以上讨论主要基于\textbf{稳态过程}。然而在实际催化过程中，催化剂的活性和选择性随反应
时间而变化：反应初期活性逐步上升（如诱导期），反应末期逐步下降（如失活）；也有人为地
周期变动操作条件（温度、浓度、组成）以提高生产能力或改善选择性。在这些非稳态过程中，
反应速率随时间而变化。另一方面，即使体系处于稳态，\textbf{稳态动力学本身也给出不了
吸附态的浓度与寿命}——它只能给出速率对流体相组成的依赖。本节的四个小节分别处理：
活性随时间衰减（积炭失活）、同一条件下的多个定态与振荡、以及两种以同位素与脉冲为工具
的瞬变方法（SSITKA 与 TAP）。\src{S2 §10.6 引言，书内 p.793 / PDF p.282}

\begin{probout}
非稳态方法的价值不在于「更高级」，而在于它测的是\textbf{另一类量}：稳态实验能给出
$r=f(c_i,T)$，却无法回答「表面上现在有多少中间物种、它们待了多久、哪一步最慢」。
本节的落点是三条时间轴——\textbf{秒以下的表面过程}（SSITKA、TAP）、\textbf{分钟到小时的
催化剂失活}（积炭动力学）、\textbf{秒到分的自持振荡}（多稳态与振荡）。选择方法的第一
原则是\textbf{时间尺度匹配}：方法的时间分辨率必须快于所测过程的时间常数，否则得到的
是卷积后的假象。
\end{probout}

\subsubsection{积炭失活动力学}

早在 20 世纪 50 年代就有关于积炭导致催化剂失活的、积炭浓度与时间的指数关系经验式，
但长期难以发展成有严格理论基础的动力学。主要原因有两层：积炭本身是一个复杂过程；
碳沉积与主反应、与扩散相互纠结，难以区分。从来源看，积炭可以是起始反应物本身的副反应
（与主反应\textbf{平行}），也可以从中间产物开始（与主反应构成\textbf{连串}）；从位置看，
碳沉积可以发生在全部活性位上，也可以选择性地发生在某些活性位上；根据扩散模数的大小或
积炭反应的不同，积炭可以发生在孔道外部或内部、颗粒外表面或核心、床层开始端或尾部；
较高积炭量时孔道会有不同程度的堵塞，堵塞程度取决于积炭量、积炭机理与孔结构。由于积炭
活性位的性质尚不清楚、孔结构也难以在几何上严格描述，严格的动力学处理很困难。Froment
等在活性位、催化剂颗粒和反应器\textbf{三个尺度}上建立了动力学方程。
\src{S2 §10.6.1，书内 p.793—794 / PDF p.282—283}

\begin{principle}
\textbf{活性位尺度的积炭失活模型。}以单分子可逆反应
$\mathrm{B_1\rightleftharpoons B_2}$ 为例。当表面反应为速率控制步骤时，按
Langmuir–Hinshelwood（LH-HW）方法，速率方程可写为
\begin{equation}
r=k_1K_1c_L\left(c_1-\frac{c_2}{K_{eq}}\right)
\label{eq:coke-r}
\end{equation}
碳的前体 C 与吸附的 $\mathrm{B_1}$、$\mathrm{B_2}$ 竞争活性位，形成不可逆的
$\mathrm{CL}$ 物种，于是总活性位浓度
\begin{equation}
c_t=c_L+c_1+c_2+c_{CL}
\label{eq:coke-sites}
\end{equation}
按 Langmuir 吸附方程，式~\eqref{eq:coke-sites} 可写成
\begin{equation}
c_t-c_{CL}=c_L(1+K_1c_1+K_2c_2)
\label{eq:coke-langmuir}
\end{equation}
把式~\eqref{eq:coke-langmuir} 代入式~\eqref{eq:coke-r}，得
\begin{equation}
r=\frac{c_t-c_{CL}}{c_t}\cdot
\frac{k_1K_1\left(c_1-c_2/K_{eq}\right)}{1+K_1c_1+K_2c_2}
=\varphi\,r_{\text{fresh}},\qquad
\varphi\equiv\frac{c_t-c_{CL}}{c_t}
\label{eq:coke-activity}
\end{equation}

\textbf{这个等式只有一种正确读法，代入时不要多做一步除法}：右边第一个因子
$\varphi=(c_t-c_{CL})/c_t$ 是\textbf{无量纲的活性系数}（$\varphi=1$ 为未失活），
第二个因子才是\textbf{速率}（单位与 $r$ 相同）。两者是\textbf{相乘}关系，
$c_t$ 只在 $\varphi$ 里出现一次。逐项核对：

- \textbf{与直接代入恒等}：式~\eqref{eq:coke-langmuir} 给出
  $c_L=(c_t-c_{CL})/(1+K_1c_1+K_2c_2)$，代回式~\eqref{eq:coke-r} 立即得到上式，
  故式~\eqref{eq:coke-activity} 与式~\eqref{eq:coke-r} 是同一个方程，增减任何因子都会破坏恒等。
- \textbf{未失活极限}：$c_{CL}=0$ 时 $\varphi=1$，$r$ 回到新鲜的
  $r_{\text{fresh}}=k_1K_1(c_1-c_2/K_{eq})/(1+K_1c_1+K_2c_2)$；$c_{CL}=c_t$
  时 $\varphi=0$、$r=0$，两端都正确。
- \textbf{对总位点浓度的依赖}：$c_{CL}$ 固定时
  $r\propto(c_t-c_{CL})$；若 $c_{CL}$ 与 $c_t$ 同比例变化（$\varphi$ 不变），
  则 $r\propto c_t$。$c_t$ 翻倍、$c_{CL}$ 不变时速率不会只翻一倍，
  而是按 $(c_t-c_{CL})$ 线性增加——这正是「活性系数」与「绝对速率」必须分开的原因，
  也是该模型给不出绝对速率、只能给出相对失活描述的根本局限。

活性位的初始积炭覆盖速率 $r_{s0}$
\textbf{不能直接测量}，只能由得到的积炭形成速率 $r_C$ 换算，换算中还要用到积炭的相对
分子质量 $m_C$ 与单个活性位上的积炭量 $\gamma$：
$r_{s0}=r_C/(m_C\gamma c_t)$。
\tocheck{未对照原页面图像：$r_{s0}$ 的换算式中 $\gamma$ 与 $c_t$ 的位置无法从 OCR 判定
（OCR 仅给出 $r_C/(c_t m_C)$）；式~\eqref{eq:coke-activity} 分子中\textbf{不应}再乘一个
$c_1$——$c_1$ 已经出现在 $(c_1-c_2/K_{eq})$ 里，再乘一次即重复计入，OCR 是否多出这一
因子未能判读。}
\src{S2 §10.6.1，书内 p.794 / PDF p.283}
\end{principle}

活性位数目的减少主要由积炭机理和生成速率决定。对于\textbf{平行}积炭反应
$r_C=c_tk_{C1}K_1c_1/(1+K_1c_1+K_2c_2)=r_{C0}\varphi_C$；对于\textbf{连串}积炭反应
$r_C=c_tk_{C2}K_2c_2/(1+K_1c_1+K_2c_2)=r_{C0}\varphi_C$。若积炭与主反应在同一活性位上
进行、并且速率决定步骤涉及相同数目的活性位，则 $\varphi_C=\varphi_1$，即积炭系数等于
主反应的失活系数。若主反应的活性位数为 $n_1$、积炭的活性位数为 $n_C$，则

\begin{equation}
\varphi_1=\left(\frac{c_t-c_{CL}}{c_t}\right)^{n_1},\qquad
\varphi_C=\left(\frac{c_t-c_{CL}}{c_t}\right)^{n_C}
\label{eq:coke-phi}
\end{equation}

式~\eqref{eq:coke-activity} 通过 $c_{CL}$ 求解 $\varphi_C$ 与 $\varphi_1$，对活性位被
覆盖而失活的情形有清晰的物理意义。该模型尽管没有考虑孔道堵塞和温度变化，对一些简单
积炭失活动力学研究仍有指导意义——例如甲烷蒸气重整中的积炭过程（积炭主要源于甲烷裂解
反应）；但对石油炼制和化工中的积炭反应要复杂得多，其动力学研究还有待深入。
\src{S2 §10.6.1，书内 p.795 / PDF p.284}

对\textbf{催化剂颗粒尺度}的积炭动力学，浓度分布与孔的网络结构必须考虑，Froment 等用
阈值理论和有效介质近似建立了相应方程；\textbf{反应器尺度}的积炭失活过程更为复杂，
可参考文献。在实用上，通常给出的是催化剂瞬时活性与积炭量的关系（比活性与时间的关系对
机理研究、反应器设计和催化剂再生更实用），已有线性式、指数式与双曲线式等不同经验式。
例如 Nam 和 Kittrell 把积炭分为单层覆盖（占据活性位）与多层覆盖（不另占活性位），并
设定多层积炭生成速率对反应物 $\mathrm{B_1}$ 的分压 $p_1$ 和单层积炭浓度 $c_\alpha$
均为一级：

\begin{equation}
\frac{\delta c_\beta}{\delta t}=k_\beta c_\alpha p_1
\label{eq:multilayer-coke}
\end{equation}

由此推出总积炭量 $c_\sum$ 与活性系数 $k_t$ 的关系

\begin{equation}
c_\sum=q_1(1-k_t)-q_2\ln k_t,\qquad
q_1=m_CS_0-q_2,\qquad
q_2=\frac{k_\beta m_CS_0(1+K_1p_1+K_2p_2)}{k_\alpha K_1}
\label{eq:nam-kittrell}
\end{equation}

式中 $c_\alpha$、$c_\beta$、$c_\sum$ 分别为单层、多层和总积炭量，$k_\alpha$ 与
$k_\beta$ 分别为单层和多层碳生成速率常数，$k_t$ 为活性系数。原书指出：在某一参数关系下
$k_t$ 与 $c_\sum$ 呈\textbf{指数关系}，在另一关系下呈\textbf{线性关系}；该模型同样未
考虑孔道堵塞与温度变化，也未考虑表面反应速率控制以外的其它机理，但对\textbf{微分反应器}
上的失活动力学数据处理非常有用。
\tocheck{未对照原页面图像：式~\eqref{eq:nam-kittrell} 中 $\ln k_t$ 项的正负号，以及
「$q_1$ 与 $q_2$ 何种相对取值给出指数关系、何种给出线性关系」——OCR 在此处丢失了两个
关系符号，无法判定。}\src{S2 §10.6.1，书内 p.795—796 / PDF p.284—285}

\textbf{非恒温条件下的失活动力学。}解析时需在物料衡算方程和失活方程中引入 Arrhenius
数，在热平衡方程中引入主反应与积炭反应的热效应。以放热反应为例，催化剂球形颗粒中
$\mathrm{B_1}$ 的浓度与温度分布可用一组\textbf{无因次方程组}描述（两个二阶方程分别给出
无因次浓度与无因次温度，源项中含失活参数与热效应参数），失活速率方程则把无因次活性
系数的时间导数写成两个并行／连串积炭函数之和；通过数值计算可求出平行积炭与连串积炭中
不同反应时间颗粒内温度与活性系数的径向分布，以及颗粒的\textbf{非恒温有效因子}。
这套分析的前提是\textbf{低积炭}，即孔道无明显堵塞；积炭量高、孔道堵塞严重影响扩散时，
动力学中必须考虑失活对扩散的影响。\tocheck{未对照原页面图像：该无因次方程组与失活速率
方程的具体形式（$\nabla^{2}$ 项的系数、失活方程右端两项的权重 $k_{f1}/k_{f2}$）在 OCR
中不完整，故此处只叙述结构与前提，不抄录方程。}
\src{S2 §10.6.1，书内 p.796 / PDF p.285}

实际催化过程中的失活动力学更为复杂，通常只能用经验式表达。例如重油催化裂化，组分非常
复杂、催化剂活性位强度分布范围很宽，即使简化处理，也会发现\textbf{积炭的反应级数从初期
的高级数向后期的低级数过渡}。此外，上述解析都假定\textbf{反应动力学本身不受失活影响}；
虽有学者从理论上论证非均匀表面上的失活动力学不宜采用这一假定，但由于表面非均匀性不易
处理、且该假定在一些失活动力学处理中应用成功，用它处理实际问题仍是可行的。
\src{S2 §10.6.1，书内 p.796 / PDF p.285}

\begin{crosscheck}
积炭失活的数据必须与\textbf{稳态动力学}分开拟合、再互查：（1）式~\eqref{eq:coke-r}—
\eqref{eq:coke-activity} 假定主反应的速率方程形式（含 $K_1$、$K_2$）在失活过程中不变，
这只在\textbf{活性位数按同一比例减少}时成立——若积炭选择性地覆盖某类活性位，
$K_1/K_2$ 也会变，此时「活性系数」不再是单一标量。（2）检验办法是把不同积炭量下的
稳态速率方程分别拟合，看 $K_i$ 是否漂移；漂移则说明失活不是简单的位点计数问题。
（3）式~\eqref{eq:nam-kittrell} 一类活性-积炭量关系式是对\textbf{微分反应器}数据最有用，
与积分床数据互查时须先做积分校正。
\end{crosscheck}

\subsubsection{催化反应动力学中的多稳态与振荡}

催化反应的稳态指体系的质量、能量流动处于平衡状态，各项参数不随时间改变。
\textbf{多稳态}的提出最初基于对强放热反应的反应器解析：在同一入口温度下，绝热反应器
可以在两个定态下操作。它在催化剂颗粒、全混流反应器、进出口自换热反应体系等不同尺度上
都存在。多稳态最初主要归因于传质过程，但在\textbf{恒温且无传质影响}的条件下仍会出现：
例如 Pt/Al$_2$O$_3$ 催化剂上的 CO 氧化，在相同恒温条件下 CO 转化率可以恒定在约 $10\%$
（「低稳态」），也可以恒定在约 $40\%$（「高稳态」），\textbf{取决于催化剂的预处理条件}。
若系统从高稳态自动转入低稳态，并随时间延长自动由低稳态转入高稳态，就在两个稳态之间
形成持续性\textbf{振荡}。例如 CO 在 Pt/Al$_2$O$_3$–SiO$_2$ 上的氧化反应中反应速率随
CO 浓度振荡；在低 CO 分压下反应对 CO 是一级，在高分压下对 CO 是\textbf{负一级}。
\src{S2 §10.6.2，书内 p.796—797 / PDF p.285—286}

振荡行为在氢气、氨气、烃类、甲醇和乙醇氧化、F–T 合成、乙烯和硝基苯加氢等反应中广泛
存在。研究它有两重意义：既揭示影响催化剂性能的因素与非稳态操作下的潜在危险，
也可以\textbf{利用振荡提高催化反应效率}（例如通过调节温度和组成强制振荡以提高转化率）。
对 Pt 上 CO 氧化振荡的研究，揭示了 Pt 表面重构、具有高 O$_2$ 黏结系数的表面缺陷、
CO 表面覆盖度变化等一系列影响催化剂性能的因素。
\src{S2 §10.6.2，书内 p.797 / PDF p.286}

\begin{principle}
\textbf{振荡的数学描述。}原理上，振荡可用总包反应中物种的表面浓度 $c_i$ 随时间变化的
一组微分方程描述，其速率方程可由基元步骤得到：
\begin{equation}
\dd c_i/\dd t=F_i(c_j,\,p_k)
\label{eq:osc-ode}
\end{equation}
式中 $c_j$ 为所有的态变量，$p_k$ 为外部参数（温度、分压等）。然而对实际体系，空间均一性
是不可能的，甚至对单晶表面也是如此。对固定床反应器，催化反应振荡是多个因素\textbf{同步
化}的结果，同步化因素包括：催化剂、载体或气相的传热；表面振荡导致的分压变化；吸附物种
的表面扩散或表面结构的转变。因此式~\eqref{eq:osc-ode} 需修正为
\begin{equation}
\partial c_i/\partial t=F_i(c_j,\,p_k)+D_i\,\frac{\partial^{2}c_i}{\partial r^{2}}
\label{eq:osc-ode-diff}
\end{equation}
式中 $D_i$ 为普适扩散系数，$r$ 为空间坐标。一般每个反应均有其自己的振荡机制，普适化
只在一定程度上可行。\src{S2 §10.6.2，书内 p.797 / PDF p.286}
\end{principle}

\begin{pitfall}
\textbf{不要把振荡区的数据当稳态数据拟合。}（1）多稳态意味着\textbf{同一操作条件对应
两个以上的速率值}，若实验不记录预处理历史与升降温路径，同一组 $(T,\,p_i)$ 会给出互相
矛盾的速率，看起来像「实验不可重复」。（2）式~\eqref{eq:osc-ode-diff} 提示：固定床中的
振荡是同步化现象，因此\textbf{反应器尺度的振荡周期与颗粒尺度的本征周期不等}，用床层
出口的振荡周期去反推表面基元步骤的时间常数是危险的。（3）工程上的正面用法是「利用」：
强制周期操作（周期进料、周期变温）可以改变时均转化率与选择性，这时应以\textbf{时均
性能}与稳态性能对比，而不是逐点比较。
\src{S2 §10.6.2，书内 p.797 / PDF p.286}
\end{pitfall}

\subsubsection{非稳态方法一：稳态同位素瞬变动力学分析（SSITKA）}

\term{SSITKA}{steady state isotopic transient kinetic analysis}{稳态同位素瞬变动力学
分析}是\textbf{同位素示踪技术}与\textbf{动态瞬变技术}的结合，可分为稳态与非稳态两类
实验：非稳态实验主要用于获得反应机理的信息，而稳态实验则融瞬变与稳态实验于一身，
在获得定性信息的同时还能获得反应中间物的\textbf{定量}信息。SSITKA 一般采用
\textbf{质谱}作为检测手段，通过检测未反应的反应物和反应产物的瞬变响应，获取反应机理和
原位动力学信息。原位红外或核磁共振在观察催化剂表面方面有独特优势，可以直接观察反应
过程中表面与表面上吸附物种的变化；把 SSITKA 与 FTIR 或 NMR 结合，利用瞬变红外或瞬变
核磁共振光谱，不仅可以研究表面动态反应和鉴别表面吸附中间物种，还可以通过\textbf{同位素
位移效应}正确表征表面中间物种——包括其化学结构和表面覆盖度，并可分辨吸附的反应物种与
非反应物种，弥补了单纯用质谱检测的不足。
\src{S2 §10.6.3，书内 p.797 / PDF p.286}

\begin{measure}
\textbf{SSITKA 的装置与操作（原书图 10-23）。}装置包括\textbf{气路控制部分}、\textbf{反应
器}和\textbf{质谱在线分析系统（MS）}，并常配一台色谱仪以测定反应物和产物的浓度。
气路控制部分最重要的是切换气路的\textbf{四通阀}：切换时应尽量不对体系的压力、气体组成
和流动产生干扰，以保持原有稳态。反应器采用\textbf{理想反应器}，如活塞流反应器或全混流
反应器；尤其是活塞流反应器体积小、操作简单，短床层具有无轴向扩散的特点。检测用的质谱仪
是 SSITKA 的核心，要求能够\textbf{多通道连续快速分析且时间分辨率高}。
\textbf{操作流程}：实验前，反应物与同位素标记的反应物两路气体需要连续流动；然后用四通阀
切换气路，使同位素标记的反应物通过反应器并在催化剂表面反应生成同位素标记的（中间）
产物，用质谱仪跟踪其\textbf{替代过程}及其变化曲线。
\src{S2 §10.6.3，书内 p.798 / PDF p.287}
\end{measure}

\begin{dataflow}
\textbf{从 SSITKA 曲线到动力学参数。}（1）对瞬变曲线积分，可获得两个基本参数：
\textbf{表面中间物种的停留时间}与\textbf{其量}。（2）若用化学吸附测定出催化剂表面暴露的
金属原子数，就可以由中间物种的量算出\textbf{表面覆盖度}。（3）若表面进行的化学反应能够
假定为\textbf{拟一级}，那么很容易求得拟一级反应速率常数。（4）SSITKA 实验还能测量与
假定反应机理相关的一些参数。利用 SSITKA 研究过的典型催化反应有：氨合成、CO 加氢、
乙烷裂解、甲烷活化和 CO 氧化等。\src{S2 §10.6.3，书内 p.798—799 / PDF p.287—288}
\end{dataflow}

\begin{pitfall}
\textbf{SSITKA 的三条使用边界。}（1）\textbf{必须有可标记元素}：它测的是同位素替代过程，
对不含可交换标记原子的反应无法施用。（2）\textbf{切换本身必须是「无扰动」的}：四通阀的
死体积、切换瞬间的压力波动与流量扰动都会直接进入瞬变曲线，把停留时间算长；这正是原书
强调「切换时应尽量不对体系压力、气体组成和流动产生干扰，保持原有稳态」的原因。
（3）\textbf{质谱的碎片与同位素峰重叠}：标记产物的母峰与碎片峰常与未标记物种的部分峰
重叠，若不做碎片校正，替代曲线会被人为拉长——常规做法是同时采集多个 $m/z$ 通道并用
纯物质标定碎片比。\tocheck{未对照原页面图像：碎片重叠与校正属常规质谱实践，原书未展开，
此条为补充，供回查。}\src{S2 §10.6.3，书内 p.798 / PDF p.287}
\end{pitfall}

\subsubsection{非稳态方法二：产物瞬态分析（TAP）}

\term{TAP}{temporal analysis of products}{产物瞬态分析}是一种研究多相催化动力学的新
技术，能够用于真空脉冲-应答实验、程序升温脱附／反应实验以及常压稳态和瞬态实验，
既可以研究模型催化剂，也可以研究工业催化剂。\src{S2 §10.6.4，书内 p.799 / PDF p.288}

\begin{measure}
\textbf{TAP 的四个组成部分与关键参数（原书图 10-24）。}（1）\textbf{高真空下的微反应器}：
可以只填充催化剂，也可以在两端填充惰性填料、中间填充催化剂——后者称为\textbf{三区
反应器}，优点是可以获得更多的信息并进一步简化数学处理；如果催化剂床层填充得非常薄，
可以认为是全混流反应器，使数学处理更简化且床层状态更均一。（2）\textbf{产生窄脉冲的
电磁脉冲发生器}：时间宽度可到 $10^{-4}$~s 数量级，弛豫时间可以在 $10^{-4}$—$10^{3}$~s
之间变化；如此短的脉冲使 TAP 能够检测到毫秒级的化学反应。（3）\textbf{质谱检测器}：
一般采用四极质谱仪检测反应物、中间产物和最终产物的化学组成；一般电离室与反应器出口间
距离不大于 $1$~cm，背景压力小于 $10^{-6}$~Pa，具有非常高的灵敏度。（4）\textbf{自动
控制／数据处理系统}。\textbf{脉冲进样量的量级}：反应物可以控制在 $10^{10}$—$10^{15}$ 个
分子，只占催化剂表面活性位数量的\textbf{百万分之一}，因此不会对催化剂的表面状态产生
影响；检测到的结果是催化剂表面恒定状态下反应分子与表面的作用，反映的是\textbf{本征的
相互作用}。又由于进样量非常少，分子对流传输消失，\textbf{仅有 Knudsen 扩散}。
\src{S2 §10.6.4，书内 p.799 / PDF p.288}
\end{measure}

\begin{dataflow}
\textbf{从 TAP 曲线到参数：标准曲线比较法。}TAP 得到的是反应物、中间产物和最终产物的
气流随时间变化的曲线；通过分析与数学处理得到动力学和传递方面的信息。目前已经发展出
处理下列情形的\textbf{理论标准曲线}：仅 Knudsen 扩散、扩散 + 不可逆反应、扩散 + 可逆
吸附、快速吸附-慢脱附、快速吸附-脱附等。把实验曲线与标准曲线比较，就可以分析判断并
测量\textbf{扩散系数、不可逆吸附／反应、可逆吸附}等参数，从而获得催化剂的物理化学和
动力学状态的信息。TAP 已用于氧化物上的吸附和脱附、氧活化、丙烯选择氧化、甲醇氧化／
氨氧化、丁烷氧化、甲烷转化等研究，给出吸附、脱附、扩散等参数以及反应网络和机理的新
论据。\src{S2 §10.6.4，书内 p.800 / PDF p.289}
\end{dataflow}

\begin{crosscheck}
\textbf{三类瞬变方法的时间尺度与信息分工，是本节最需要记住的一张表。}
TAP 的脉冲宽度在 $10^{-4}$~s 量级、弛豫时间 $10^{-4}$—$10^{3}$~s，因此它能分辨
毫秒级过程且几乎不扰动表面；SSITKA 的典型时间窗是表面中间物种的停留时间（秒量级），
且必须在\textbf{稳态背景}下做同位素切换；过渡应答法施加的是浓度阶跃，其应答衰减在
秒到分量级；而积炭失活与振荡分别是分钟—小时与秒—分尺度。\textbf{方法能不能用，首先
取决于它比所测过程快多少}。三者与稳态方法的分工见表~\ref{tab:ch10-transient}。
\end{crosscheck}

\begin{table}[tbp]
\centering
\caption{稳态与非稳态动力学方法的时间尺度与信息分工。看第二列：\textbf{方法的时间分辨率
必须快于所测过程的时间常数}，否则得到的是与仪器应答卷积后的假象；看第三、四列：
「给不出」一栏决定了任何单一方法都不能独立定出机理。\src{据 S2 §10.5.3—§10.5.4、
§10.6.1—§10.6.4，书内 p.788—800 / PDF p.277—289 整理。}}
\label{tab:ch10-transient}
{\srcsize
\begin{tabularx}{\twocolwidth}{@{}p{21mm}p{26mm}p{52mm}X@{}}
\toprule
方法 & 特征时间尺度 & 能给出什么 & 给不出什么 \\
\midrule
稳态动力学 & 达稳态后取数（分钟—小时） & 速率方程、速率常数、表观活化能、反应级数、
吸附平衡常数的组合、选择性的宏观规律 & 吸附态浓度与寿命、中间物种的身份、速率控制步骤的
直接证据；无法区分「$K$ 大而覆盖度高」与「$K$ 小而覆盖度低」 \\
过渡应答（浓度阶跃） & 秒—分（应答衰减） & 吸附量（应答曲线下的阴影面积）、吸附与
脱附速率、由产物应答的即时性判断速率控制步骤、由降零实验判断基元步骤可逆性 & 绝对覆盖度
（需先测活性位数）；解析依赖假定的基元步骤网络，网络错则参数系统性偏移 \\
SSITKA & 秒量级（表面中间物种停留时间） & 中间物种的\textbf{停留时间与量}；配化学吸附
得覆盖度；拟一级假定下得拟一级速率常数；与原位 IR／NMR 联用可得中间物种结构与同位素
位移信息 & 毫秒以下过程；必须存在可标记元素；切换瞬间的压力与流量扰动会直接污染停留
时间 \\
TAP & 脉冲 $10^{-4}$~s；弛豫 $10^{-4}$—$10^{3}$~s & 仅 Knudsen 扩散下的扩散系数、
不可逆吸附／反应、可逆吸附参数；反应网络与机理新论据；几乎不扰动表面（脉冲分子数只占
表面活性位的百万分之一） & 常压稳态速率与高转化率选择性；对流传输消失，故不能代表工业
条件下的相间传递 \\
程序升温脱附／反应（TAP 的升温模式，即 TPSR 类实验） & 由升温速率决定，通常为分钟级
& 表面吸附态与中间物\textbf{随温度的转化序列}，以及产物与活性位的对应关系 & 单独的定量
动力学参数——温度与覆盖度同时变化，必须借助模型反演，或与 SSITKA／稳态数据联合定参
\tocheck{未对照原页面图像：原书只在 TAP 的能力列表中列出「程序升温脱附／反应实验」，
未给出升温速率范围；本行的分钟级时间尺度与信息表述属常规实践归纳，需回查原书或另找
出处。} \\
积炭失活 & 分钟—小时—天 & 活性系数随积炭量的衰减规律（线性／指数／双曲线型）；三段
尺度（活性位、颗粒、反应器）的失活模型 & 初始积炭覆盖速率 $r_{s0}$ 不能直接测，只能由
$r_C$ 换算；高积炭下的孔道堵塞须另行处理 \\
多稳态与振荡 & 秒—分（振荡周期） & 同一操作条件下的多个定态与自持振荡；表面重构、
覆盖度变化等机理线索；强制周期操作可改变时均性能 & 用出口振荡周期反推表面基元步骤的
时间常数（固定床中是同步化结果）；也不能当作稳态速率数据使用 \\
\bottomrule
\end{tabularx}}
\end{table}

\subsection{结论}

多相催化动力学已发展成为内容非常丰富的交叉学科，近五十年来在\textbf{宏观-介观-微观}
不同尺度上都有不同程度的发展。宏观尺度上的动力学是催化剂性能最直接的表征手段，
但其所能表征的性能与反应的复杂程度直接相关：对\textbf{简单反应}，速率方程、速率常数、
活化能、反应级数及吸附平衡常数等可以较准确而全面地表征；对\textbf{复杂反应}，可通过
反应网络的解析求出催化剂上初级和次级反应的速率常数，从而在更深层次认识催化反应的本质。
在多相催化反应中，\textbf{传递过程也是催化剂的一个重要性能}，对实际催化过程尤其如此。
宏观多相催化动力学还为催化过程提供了数学模型：由于实验技术、分析方法和数值模拟的发展，
所得模型基本可以准确地反映温度、压力和空速等工艺参数对催化反应速率、转化率和选择性的
影响规律，是催化反应器设计和过程优化的基础与手段。此外，宏观多相催化动力学也有助于
催化剂设计及性能优化，以及催化反应机理等基础认识的发展。
\src{S2 §10.7，书内 p.800—801 / PDF p.289—290}

\begin{crosscheck}
把本章三条线索接起来：\textbf{分片 a} 给出「RDS + 吸附平衡 → 速率方程」的形式框架；
\textbf{分片 b} 给出「传递是否限制」的判据与有效因子；\textbf{本分片}给出「怎样用反应器
和实验把二者分开测量」，以及稳态方法给不出的那部分信息如何由非稳态方法补上。任何一条
速率方程在报告时都应同时交代三件事：\textbf{在哪个动力学区取得}（本征／内扩散／外扩散）、
\textbf{传递判据是否通过}（$Da$、$\Phi$、Mears 判据）、\textbf{吸附参数是否独立验证过}
（过渡应答、SSITKA 或 TAP）。缺任何一项，方程就只能作为经验拟合式使用。
\end{crosscheck}

\srcnote{本分片覆盖 S2（下册）第 10 章 §10.5（动力学测定方法和实验装置，含
§10.5.1 实验室反应器、§10.5.2 内扩散参数的测定、§10.5.3 过渡应答法、§10.5.4 本征
动力学区的确定和排除传递过程干扰）、§10.6（非稳态催化过程动力学，含 §10.6.1 积炭失活
反应动力学、§10.6.2 多稳态与振荡、§10.6.3 SSITKA、§10.6.4 TAP）与 §10.7 结论，
源页面为书内 p.779—806 / PDF p.268—295。表~\ref{tab:ch10-reactor} 据原书表 10-3 与
§10.5.1 正文整理，其中「浆态（液固）」一行为补充；图~\ref{fig:ch10-velocity} 与
图~\ref{fig:ch10-etaphi} 为模型计算图，非原书插图复制。公式的重建与存疑点已在正文中
以 \tocheck{} 逐处标明。}
